% ============================================================================
%  CAPÍTULO II: ALCANCE Y LÍMITES
% ============================================================================

\section{Alcance y Límites}

\subsection{Alcance}

El presente informe documenta la configuración de los servicios DHCP y web
del proyecto. Las actividades realizadas incluyen:

\begin{itemize}
  \item Instalación y configuración del servidor DHCP Kea sobre Debian, con
        definición de la subred \texttt{192.168.0.0/24}, el ámbito de
        asignación dinámica para los equipos de usuario
        (\texttt{192.168.0.110}--\texttt{192.168.0.254}) y las opciones de
        red entregadas a los clientes (gateway, servidor DNS y dominio).
  \item Desactivación del servicio DHCP del router de la red para evitar
        conflictos de asignación entre servidores.
  \item Verificación de la asignación automática de direcciones IP desde
        equipos clientes y del registro de concesiones en el servidor.
  \item Instalación y configuración nativa de Nginx sobre Debian GNU/Linux
        mediante el gestor de paquetes APT.
  \item Configuración de virtual hosts para el sitio principal del proyecto
        (\texttt{web.sudoers.lan}, \texttt{www.sudoers.lan},
        \texttt{sudoers.lan}).
  \item Configuración de proxy inverso hacia los servicios internos:
        \begin{itemize}
          \item CUPS (impresión): \texttt{print.sudoers.lan}
          \item Filebrowser (archivos): \texttt{files.sudoers.lan}
          \item Portainer (panel de administración): \texttt{portainer.sudoers.lan}
          \item Zabbix (monitoreo): \texttt{zabbix.sudoers.lan}
        \end{itemize}
  \item Generación de certificados SSL/TLS auto-firmados mediante la herramienta
        OpenSSL con directivas SAN (\textit{Subject Alternative Name}).
  \item Configuración de redirección HTTP a HTTPS en todos los virtual hosts.
  \item Configuración del servidor por defecto (\texttt{default\_server}) para
        capturar peticiones no mapeadas.
  \item Configuración de soporte WebSocket para Portainer.
  \item Despliegue de la landing page institucional como contenido estático en
        \texttt{/var/www/html}.
  \item Habilitación y control de los demonios Nginx y Kea a través del
        sistema de inicio \texttt{systemd}.
  \item Pruebas de acceso desde navegadores, verificación de los proxies
        inversos y comprobación de la asignación DHCP desde clientes.
\end{itemize}

\subsection{Límites}

Los siguientes aspectos no forman parte de este informe:

\begin{itemize}
  \item Certificados SSL emitidos por una autoridad de certificación (CA)
        confiable; se utilizan certificados auto-firmados adecuados para
        un entorno de laboratorio.
  \item Integración con dominios externos o configuración de registros DNS
        públicos.
  \item Balanceo de carga o alta disponibilidad del servidor web o del
        servidor DHCP.
  \item Configuración de WAF (\textit{Web Application Firewall}) o reglas
        avanzadas de seguridad en Nginx.
  \item Seguridad avanzada como protección contra ataques DDoS, rate limiting
        o listas de control de acceso por IP.
  \item Configuración de servidores DHCP redundantes (\textit{failover}) o
        sincronización de concesiones entre múltiples nodos.
  \item DHCPv6 y la gestión de rangos de direcciones IPv6.
  \item Reservas estáticas de direcciones por MAC en el servidor DHCP.
\end{itemize}

\subsection{Entorno de trabajo}

El entorno utilizado para la configuración consta de los siguientes
componentes:

\begin{itemize}
  \item Servidor: máquina física con Debian GNU/Linux 12 (Bookworm) donde se
        instalan y configuran los servicios de forma nativa.
  \item Despliegue: administración directa mediante paquetes oficiales de
        Debian y orquestación de servicios a través de \texttt{systemd}.
  \item Red interna: red administrativa \texttt{192.168.0.0/24}
        (WLO1) y red de usuarios \texttt{192.168.20.0/24} (ENO1).
  \item Servicios: Nginx (Web), Kea DHCP (Debian GNU/Linux 12 Bookworm).
  \item Resolución de nombres: servidor DNS propio (BIND9) que resuelve
        el dominio \texttt{sudoers.lan} para todos los servicios.
  \item Clientes de prueba: máquinas Linux y Windows conectadas a la
        red local, navegadores web para verificación de los sitios y
        proxies inversos, y clientes DHCP para la comprobación de la
        asignación automática de direcciones.
  \item Usuarios: \texttt{nicolas}, \texttt{joaquin} y
        \texttt{david}, integrantes del equipo \textit{Grupo 1 Sudoers}.
\end{itemize}